\documentclass[11pt]{article}
\usepackage[T1]{fontenc}
\usepackage[utf8]{inputenc}
\usepackage{lmodern}
\usepackage[margin=1in]{geometry}
\usepackage{amsmath,amssymb,mathtools,mathrsfs}
\usepackage{microtype}
\usepackage{xcolor}
\usepackage{hyperref}
\hypersetup{colorlinks=true,linkcolor=black,urlcolor=blue,citecolor=black,pdftitle={Semantic Distinction}}
\usepackage{bookmark}
\usepackage{fancyhdr}
\pagestyle{fancy}\fancyhf{}\fancyhead[L]{\small Semantic Distinction}\fancyfoot[C]{\thepage}
\setlength{\headheight}{14pt}
\setlength{\emergencystretch}{3em}
\setlength{\parskip}{0.25em}
\providecommand{\tightlist}{\setlength{\itemsep}{0pt}\setlength{\parskip}{0pt}}
\title{Semantic Distinction\\\large Behavioral Consequences and Contextual Neural Differentiation}
\author{Flyxion}\date{September 2026}
\begin{document}
\maketitle
\begin{abstract}
Franch and colleagues report a positive correspondence between
hippocampal population-response distances and contextual language-model
embedding distances that changes sign within a restricted range of
closely related embeddings. This geometrical observation does not
establish an adaptive contrastive mechanism, and neural--model alignment
does not establish algorithmic identity. This commentary advances a
narrower hypothesis: independently normed, probe-relative consequences
of substituting situated interpretations predict neural differentiation
beyond embedding, linguistic, and fitted-mapping baselines. A controlled
experiment would measure substitution consequences in a separate
behavioral sample before neural analysis, then test their association
with population distances and held-out discrimination. The proposal
distinguishes representational geometry from recoverable information and
task consequence, specifying which observations would support contextual
differentiation while leaving its biological mechanism open.

\end{abstract}
\noindent\textit{Theoretical commentary. No new experimental data or fitted analyses are reported.}

\hypertarget{the-problem-of-preserving-relations}{%
\section{The problem of preserving
relations}\label{the-problem-of-preserving-relations}}

Understanding an expression requires recognizing what it shares with
other expressions while retaining what makes its present occurrence
different. A representation that supports only similarity risks
conflating distinctions; a representation that isolates every occurrence
loses the regularities through which comprehension becomes possible.
Semantic organization consequently involves a tension between
generalization and discrimination. This tension becomes especially
visible when related expressions occupy nearby positions in an embedding
space, yet their substitution changes the interpretation of a sentence.
A category can remain stable while the relation asserted within it
changes. The difference between an object being available and being
unavailable, for example, is small in the vocabulary required to
describe it and substantial in the continuation that the description
permits.

The hippocampal recordings reported by Franch et al.~(2026) provide an
empirical occasion for examining this problem. Their study concerns
neural responses during narrative listening, and its principal
contribution is evidence about population representation rather than a
complete theory of meaning. The present paper develops an interpretation
of that evidence: semantic geometry may preserve relations through an
interaction between contextual integration and selective
differentiation. This interpretation is a theoretical extension, not a
claim that the experiment directly measured the practical consequences
of semantic distinctions.

The explanatory object is contextual discriminability within population
representations. Relational memory supplies background motivation for
how occurrences might be bound to discourse, but the proposed experiment
does not identify binding by measuring separation alone. ``Relational''
is therefore used below for the relevant memory literature or for
geometrical relations among responses, with those senses distinguished
explicitly. The central claim concerns whether a population preserves
distinctions required by an independently specified comprehension task.
It neither treats relations as an exhaustive account of semantics nor
dispenses with reference and the bodily or social conditions of
understanding.

\hypertarget{predictive-compression-as-a-theoretical-perspective}{%
\subsection{Predictive compression as a theoretical
perspective}\label{predictive-compression-as-a-theoretical-perspective}}

Tishby, Pereira, and Bialek's information bottleneck method (1999;
archived 2000) formalizes compression of an input X while preserving
information about a relevance variable Y. In schematic form, a
representation T is selected by minimizing \(I(X;T)\)-beta I(T;Y), with
mutual information I and trade-off parameter beta. Applied here as a
theoretical perspective, X could be the speech history and Y an
independently specified future comprehension outcome. Surface
differences may then be compressed when they carry no information about
Y, while small differences remain relevant when they change its
conditional distribution. Neither lexical identity nor acoustic
variation is intrinsically irrelevant; relevance depends on Y. This
framework motivates the consequence measure without identifying it with
mutual information.

The application is conditional rather than a demonstrated optimization
process in hippocampus. Compression and predictive retention alone do
not determine a Euclidean or cosine geometry, a local correlation sign
change, or an evolutionary explanation. Even information-equivalent
representations can have different distances under invertible changes of
coordinates. A bottleneck account of the reported geometry would
additionally require a specified representational family, noise model,
relevance variable, and readout. The useful connection to Section~5 is
that preserving similarity must be evaluated relative to what a
representation needs to predict; it is not a derivation of contrastive
separation from efficient coding alone.

For a discrete illustration, mutual information is

\begin{equation}
I(U;V)=\sum_{u,v}p(u,v)\log\frac{p(u,v)}{p(u)p(v)},\qquad
\mathcal J_{\mathrm{IB}}=I(X;T)-\beta I(T;Y).
\end{equation}

The sum uses the convention that zero-probability terms contribute zero.
Compression concerns dependence between input and representation;
predictive retention concerns dependence between representation and the
declared relevance variable. If two contexts induce the same conditional
distribution of Y, merging their codes need not lose predictive
information about that Y. If their conditional distributions differ,
merging may incur a predictive loss. This supplies a rationale for
task-relative consequentiality, without equating the finite probe score
Q with an information-theoretic optimum. In continuous deterministic
representations, mutual-information expressions can be infinite or
require additional regularity; applying this objective to real-valued
neural codes would require an explicit stochastic encoder or another
declared regularization. The commentary therefore uses the framework
conceptually rather than assigning an unmeasured information rate to the
hippocampus.

\hypertarget{what-the-hippocampal-study-establishes}{%
\section{What the hippocampal study
establishes}\label{what-the-hippocampal-study-establishes}}

Franch et al.~recorded 356 hippocampal neurons in ten native
English-speaking patients undergoing monitoring for epilepsy.
Participants listened to approximately 47 minutes of narrative speech
containing 7,346 word occurrences and 1,351 unique words. Encoding
models related neural activity to linguistic features, including
semantic embeddings, with analyses controlling for phonetic and
grammatical information. Individual neurons exhibited mixed selectivity
across words and semantic categories, while population patterns carried
information about contextual semantic organization. These observations
support distributed representation under the conditions examined,
without requiring every neuron to contribute equally or eliminating more
selective responses elsewhere in the brain.

The stimulus consisted of six Moth podcast monologues. Firing-rate
epochs began 80 ms after word onset and lasted for that word's duration,
with rates normalized by duration. The Word2Vec Poisson ridge encoding
model fit 222 of 356 recorded neurons, or 62\%; its median explained
variance was 26\%, including duration and interaction contributions. The
article describes isolated single neurons in its Results, while Extended
Data Fig. 1 illustrates single- and multi-unit isolation; a reanalysis
should verify the released unit classifications rather than infer their
proportions from this illustration. The comparison uses the final BERT
representation designated layer 13 and GPT-2 layer 37 in the article's
indexing convention.

The study also compares distances among neural responses with distances
among language-model embeddings. Contextual representations from BERT
and later GPT-2 layers corresponded more closely to the observed
population organization than static Word2Vec representations. Across the
full comparison, neural and contextual semantic distances showed small
positive correlations. Within a restricted range of especially similar
embeddings, the relationship became negative, a pattern the authors
interpret as possible contrastive coding. They report that the local
reversal was associated with different word pairs rather than repeated
occurrences of identical words and that related patterns appeared in
analyses conducted separately for individual patients.

The numerical strength and the inferential character of these findings
both matter. Overall distance correlations were approximately 0.06 for
BERT and 0.09 for GPT-2. The authors estimate that these values capture
substantial fractions of the recoverable association under their noise
model, but a noise-ceiling estimate is conditional on the assumptions
used to construct it. Encoding-model performance also includes
contributions from duration and interactions, so total explained
variance should not be described as variance attributable exclusively to
semantics. Population decoding of semantic embeddings explained
approximately 1--6\% of variance, depending on the embedding model.
These results establish measurable correspondence; they do not establish
reconstruction of a listener's complete meaning representation.

For exact numerical anchoring, Fig. 3d reports BERT r=0.0635, and
Extended Data Fig. 7a reports GPT-2 layer-37 r=0.0932; the Results prose
gives nearby values of 0.0650 and 0.0921, respectively. The
caption-specific values should not be silently treated as identical to
the prose values. The authors report approximately 54\% and 76\% of
their modeled noise ceilings for BERT and GPT-2. Those percentages
describe modeled recoverability, not percentages of semantic content
decoded or a proof of a shared mechanism. Fig. 3f reports r=-0.1536 in
the BERT distance interval 0.15--0.35, containing 2.16\% of the pairwise
observations. This restricted coverage increases sensitivity to range
boundaries and composition of the eligible pairs; the pair count does
not supply an equivalent number of independent biological observations.
The corresponding negative association for different-word pairs is
reported as approximately -0.151, while identical-word pairs show a
small nonsignificant positive association. Reduced-rank decoding reaches
patient-averaged 1-NSE=0.06 for GPT-2, with a median optimal rank of 12;
neither number should be read as a complete dimensionality estimate of
hippocampal semantics.

Frequency and polysemy also require distinctions between analyses.
High-frequency categories are centrally located in the two-dimensional
projection while showing greater separation in the full population
space. In the regression predicting distance across occurrences of the
same word, polysemy and duration are significant predictors. The
separate pair-distance regression includes semantic distance, frequency,
duration, category, and polysemy differences. These results concern
particular distance outcomes, not an unrestricted claim that polysemy
and duration explain all neural variance.

The authors explicitly acknowledge that controlled follow-up studies are
needed to test their contrastive interpretation. That acknowledgment
sets the appropriate evidential boundary for this commentary. The
distance reversal is an observation; a mechanism that increases
separation to reduce semantic confusion is an explanatory hypothesis.
The two deserve different verbs and different standards of confirmation.

The local negative association is a departure from the modest positive
associations just reported, rather than a reversal of an otherwise
near-perfect correspondence. Its interpretation depends on the
definition of the local range. The article reports approximately
0.15--0.35 for a focal BERT analysis and approximately 0.14--0.4 for
related analyses; these numerical cosine-distance ranges refer to the
particular comparison embeddings and do not define a universal threshold
of semantic similarity. A selected range can emphasize curvature or a
local fluctuation, especially when its location and endpoints are chosen
after viewing the association. The commentary therefore treats the
reported sign change as a candidate scale-dependent pattern whose
stability should be tested across independently specified ranges, rather
than as an established boundary between two coding mechanisms.

\hypertarget{relational-representation-in-the-existing-literature}{%
\section{Relational representation in the existing
literature}\label{relational-representation-in-the-existing-literature}}

\hypertarget{relational-memory-and-language-use}{%
\subsection{Relational memory and language
use}\label{relational-memory-and-language-use}}

Konkel and Cohen (2009) characterize hippocampal memory through binding
and flexible retrieval of spatial, sequential, and associative
relations; Duff and Brown-Schmidt (2012) extend relational binding and
representational flexibility to discourse and shared knowledge. These
accounts motivate investigating contextual integration during language
use without treating distances between word-evoked responses as
measurements of relational binding. Geometry among responses and memory
for experienced relations remain connected hypotheses rather than
interchangeable observations. The present experiment tests contextual
discriminability; identifying a binding contribution would require a
separate relational manipulation.

\hypertarget{separation-and-integration-as-complementary-demands}{%
\subsection{Separation and integration as complementary
demands}\label{separation-and-integration-as-complementary-demands}}

McClelland, McNaughton, and O'Reilly (1995) formulate complementary
learning systems as an account of how rapid acquisition of specific
experiences can coexist with the gradual extraction of shared structure.
Their computational argument identifies interference as a problem for
learning systems and motivates different representational and learning
properties in hippocampus and neocortex. This provides a precedent for
the tension developed here, although the present question concerns
differentiation during contextual language processing rather than the
complete architecture of consolidation. The commentary does not
introduce the general problem of preserving specificity while learning
regularities.

Bakker et al.~(2008) use high-resolution functional imaging to examine
responses consistent with pattern separation in the combined human
CA3/dentate-gyrus region. Their evidence links similar inputs to
hippocampal processing that distinguishes them, but its measurement
scale and stimulus paradigm differ from single-neuron responses during
narrative listening. Pattern separation is conventionally understood as
a transformation of input similarity into less overlapping output
representations. A negative association between neural distances and
distances in a comparison embedding is not by itself a direct
measurement of that transformation, since the embedding is not a
recording of the hippocampus's biological input. Anatomical localization
also matters: a hippocampal population result does not alone establish a
dentate-gyrus mechanism.

Schlichting, Mumford, and Preston (2015) provide a particularly close
precedent for treating integration and differentiation jointly. Their
study of overlapping associations identifies learning-related
integration and separation signatures in different hippocampal and
prefrontal regions, with learning conditions influencing the outcome.
The result cautions against assigning a single representational policy
to the entire hippocampus. It also motivates a distinction between the
separation of events that share elements and the integration of elements
that connect events. The proposed semantic account extends this question
to occurrences in discourse; it does not claim that a word-pair distance
reversal identifies the same process or reproduces the same regional
dissociation.

Together, these precedents constrain the interpretation of
``contrastive.'' It can name a descriptive geometrical pattern, a
hypothesized input--output transformation, or an adaptive learning
process that changes representations over time. These meanings require
different evidence. The Franch et al.~result supplies a candidate
descriptive pattern; the experiment proposed below would test its
sensitivity to contextual demands. Establishing a learning mechanism
would additionally require observations of representational change under
controlled exposure or intervention.

\hypertarget{mixed-selectivity-and-population-readout}{%
\subsection{Mixed selectivity and population
readout}\label{mixed-selectivity-and-population-readout}}

Rigotti et al.~(2013) show that nonlinear mixed selectivity in monkey
prefrontal cortex can support high-dimensional representations and
flexible readout during a sequence-memory task. Their analyses connect
population dimensionality to behavioral performance. This establishes a
computational precedent for information residing in combinations of
responses that are not easily interpreted neuron by neuron. It does not
make mixed selectivity a sufficient condition for semantic
representation, and selectivity across several word categories should
not automatically be equated with the particular nonlinear interactions
analyzed in their task. The hippocampal findings motivate a population
account, while additional analyses would be needed to establish which
interaction structure supports contextual discrimination.

Mixed selectivity thus motivates a population-level inquiry, while its
functional interpretation requires the noise-aware readout and
independent behavioral criterion specified in Section~5.2.

\hypertarget{contextual-meaning-within-a-wider-semantic-architecture}{%
\subsection{Contextual meaning within a wider semantic
architecture}\label{contextual-meaning-within-a-wider-semantic-architecture}}

Deniz et al.~(2023) show that linguistic context affects semantic
representations and their measurement across words, sentences, and
narratives, motivating the situated-occurrence analysis without equating
every context effect with polysemy. Patterson, Nestor, and Rogers (2007)
locate conceptual knowledge within a distributed architecture involving
modality-specific representations and a transmodal anterior temporal
hub. These precedents constrain the hippocampal interpretation:
contextual differentiation can coexist with relatively enduring
conceptual structure, and hippocampal coding does not make that region
the sole repository of meaning. Neither architecture nor contextual
variability alone predicts the local distance sign change; the
operational question is which independently consequential distinctions
remain recoverable.

\hypertarget{language-model-correspondence-and-its-evidential-level}{%
\subsection{Language-model correspondence and its evidential
level}\label{language-model-correspondence-and-its-evidential-level}}

Goldstein et al.~(2022) report correspondences between neural responses
during narrative comprehension and autoregressive language models,
including contextual representation, anticipatory prediction, and
responses associated with surprise. Their work establishes precedent for
using naturalistic language-model predictions to investigate neural
processing. These are specified computational correspondences under the
analyses performed, not evidence that every operation or learning
process is shared. The hippocampal result should be interpreted within
that comparative program while preserving its distinct anatomical target
and measurement scale.

Goldstein et al.~(2024) examine contextual embedding geometry using
dense intracranial recordings from the inferior frontal gyrus of three
participants listening to a podcast. Their zero-shot mapping analyses
relate neural representations of held-out words to language-model
geometry and find an advantage for contextual over static embeddings.
This provides a direct precedent for studying common geometrical
patterns in language representations. Franch et al.~extend the
comparison to hippocampal single-neuron populations and identify a
scale-dependent discrepancy. The difference is theoretically
informative, but the regional and recording differences prevent treating
it as a controlled comparison of hippocampal and frontal coding
policies.

The explicit puzzle is that contextual embeddings outperform static
representations in both studies, whereas the hippocampal paper reports a
local sign change that is not established as a matched frontal result in
this comparison. Is this discrepancy related to hippocampal
differentiation, or would equivalent sampling and local-distance
analyses reveal it in frontal responses as well? Lack of a matched
report is not evidence of absence. A regional claim requires the same
stimuli, comparable measurement and preprocessing, and prespecified
range analyses in both regions.

The contribution of the present commentary is correspondingly narrower
than the discovery of vectorial or contextual semantics. It proposes
that independently measured consequences of confusion could help explain
when geometrical separation becomes useful. Neither shared geometry nor
disagreement establishes that explanatory variable; it must be
manipulated or measured outside the embedding and neural-distance
analyses.

\hypertarget{representational-similarity-and-the-limits-of-geometrical-inference}{%
\subsection{Representational similarity and the limits of geometrical
inference}\label{representational-similarity-and-the-limits-of-geometrical-inference}}

Kriegeskorte, Mur, and Bandettini (2008) introduce representational
similarity analysis as a framework for comparing relations among
activity patterns with computational and behavioral models. Its
relevance here is that correspondence among distance structures concerns
relations within representations, even when individual measurement
channels have no direct counterparts. They also emphasize that distances
sharing observations are dependent and motivate inference through
condition-label randomization. This precedent supports treating geometry
as an object of inquiry while requiring uncertainty estimates that
respect the structure of the observations. It also makes independent
behavioral comparison central to the proposed functional interpretation.

A related methodological point follows from considering a linear mapping
\(r=Az\) from contextual embeddings z to neural responses r. Even
without an additional term, a general A need not preserve cosine
distance or local distance order. A geometrical discrepancy can
therefore occur without an additional separation mechanism. Evidence for
selective differentiation requires comparison with mappings that already
allow such distortions, followed by evaluation on held-out conditions.
Otherwise, an unexplained residual risks naming whatever the initial
mapping fails to capture. The strongest version of the proposal asks
whether contextual confusion consequences explain reproducible structure
beyond suitable contextual, linguistic, and mapping baselines.

This focused narrative review identifies precedents and evidential
boundaries rather than claiming exhaustive coverage. Its synthesis is
that relational binding, mixed selectivity, context dependence, and the
integration--separation tension are established research themes. The
open question advanced here concerns their conjunction in semantic
discrimination: whether the practical consequence of confusing situated
occurrences predicts population differentiation beyond the comparison
embedding distance.

\hypertarget{coordinates-metric-structure-topology-and-readout}{%
\subsection{Coordinates, metric structure, topology, and
readout}\label{coordinates-metric-structure-topology-and-readout}}

Recent work gives empirical motivation for separating coordinate-level
tuning from population organization. Yan et al.~(2026) report shared
semantic geometry across English and Spanish in human hippocampal
recordings despite differences in individual-neuron tuning, with
language-specific readout axes within a common recorded population.
Their result supports distinguishing metric correspondence from
coordinate identity. The readout axes identified analytically do not
themselves establish which downstream biological circuit decodes those
responses, and the authors' interference interpretation remains a
functional hypothesis. Shared cross-language geometry also does not
demonstrate the consequence-dependent differentiation proposed here.

Topology supplies a complementary description rather than another name
for distance. Brown and Farivar (2025) use persistent homology to
characterize components, loops, and higher-dimensional structure across
distance thresholds, showing why high correlation between dissimilarity
matrices can conceal differences in shape. This does not mean that a
complete distance matrix contains no topological information: their
analysis itself uses distance relations. The limitation concerns the
scalar comparison applied to those matrices. Coordinates, the declared
metric, topology of the sampled or hypothesized support, and
recoverability under a readout should therefore be distinguished as
structures with different invariances, rather than treated as a strict
ladder of explanations. None alone supplies an independent measure of
comprehension consequence.

\hypertarget{situated-occurrences-as-representational-units}{%
\section{Situated occurrences as representational
units}\label{situated-occurrences-as-representational-units}}

A lexical type provides a convenient index for grouping occurrences, but
grouping is not an explanation of how those occurrences are understood.
The same expression can acquire different relations through its position
in a narrative, its grammatical environment, and the expectations
established by preceding speech. Franch et al.~examine this contextual
variability through both embedding-based measures and a WordNet measure
of polysemy. Their findings associate variation across occurrences with
variation in neural responses. The distinction between an
embedding-based contextual measure and a lexical inventory of senses
remains important, since contextual variability can occur without a
clean division into dictionary senses.

The theoretical consequence is a change in the unit of analysis. Instead
of beginning with a representation r(w) assigned once to word w,
consider a representation r(w,c), where c denotes the context relevant
to that occurrence. This notation does not explain how context is
computed, how much preceding material is retained, or whether different
neural populations contribute distinct contextual variables. It merely
prevents lexical identity from being mistaken for representational
identity. Two occurrences of the same word may remain recognizably
related while differing in their encoded contribution to the current
discourse.

More precisely, write the occurrence as
\(o_i\)=(\(w_i\),\(t_i\),\(e_i\)), identifying its lexical type, onset
time, and narrative episode, and write \(c_i\)=C(\(H_{t_i}\),\(e_i\)),
where H denotes the linguistic history available before or during the
measured response. Context is not assumed to be the entire recording
session, nor is it restricted in advance to one sentence. Candidate
operationalizations should distinguish preceding words, sentence-level
information, and the currently maintained narrative episode, with
separate variables for prior knowledge and experimentally supplied task
instructions. The present commentary brackets the biological computation
of C and the duration of its effective memory; it proposes comparing
these candidate scopes. Model context windows and biological integration
timescales must also remain distinct, since a model window is an
analytical choice rather than evidence that a listener retains exactly
that interval. Within each tested condition, \(c_i\) should be specified
from the stimulus and design independently of the neural outcome.

Context sensitivity also changes the interpretation of averaging. An
average over occurrences can reveal a stable lexical component, yet it
can obscure systematic contextual distinctions. Conversely, variation
among occurrences need not all be meaningful: recording noise,
adaptation, timing, and unmeasured narrative variables can contribute. A
relational analysis should therefore seek reproducible context-dependent
structure within the variation, rather than treating either the average
or the variability as automatically authoritative.

\hypertarget{mathematical-objects-and-distinction-preservation}{%
\subsection{Mathematical objects and distinction
preservation}\label{mathematical-objects-and-distinction-preservation}}

The operational hypothesis requires interpretation, comparison
representation, response variability, and task outcome to be specified
separately, so the following vocabulary fixes their domains before any
distance or consequence is compared. Let \(\mathcal O\) be the set of
situated occurrence records, \(\mathcal K\) the experimental conditions,
\(\mathcal M\) a declared interpretation space, \(\mathcal Z\) the
comparison-representation space, and \(\mathcal Y\) the task-outcome
space. An occurrence record \(o_i=(w_i,t_i,e_i)\) identifies a token
through its lexical type, onset, and episode; its analytical context
\(c_i^{(k)}\) is supplied by the history and design under condition
\(k\). Distinct occurrence records may share a lexical type. For the
proposed experiment, \(\mathcal M\) is a finite set of explicitly
specified discourse interpretations admissible for the stimulus, such as
alternative assignments of a referent or relation. Each interpretation
is an experimental description with declared propositions and
probe-answer implications, rather than an identified biological state or
a complete theory of meaning. \(P(M\mid o_i,k)\) records uncertainty
over those alternatives, including an indeterminate category when the
protocol permits it. Equality of interpretation labels, selection of an
admissible alternative, and a protocol-defined substitution are
permitted operations; no intrinsic metric on \(\mathcal M\) is assumed.
The operational measurements of \(Y\) and \(Q\) are more developed than
this interpretation space, intentionally, because the hypothesis does
not require exhaustive semantic states. Neither an interpretation label
nor lexical identity determines a unique neural response. Lexical
non-identification is the relevant principle: shared lexical type does
not determine shared context, interpretation distribution, or population
response.

The comparison encoder and neural measurement have different domains and
noise commitments:

\begin{equation}
E:\mathcal O\times\mathcal K\longrightarrow\mathcal Z,\qquad
\mathscr R_s:\mathcal O\times\mathcal K
\longrightarrow\mathcal P(\mathbb R^{n_s}),\qquad
r_{si}^{(k)}\sim\mathscr R_s(o_i,k).
\end{equation}

Here \(\mathcal P(\mathbb R^{n_s})\) denotes response distributions for
participant \(s\), while \(z_i^{(k)}=E(o_i,k)\) is the model
representation under declared preprocessing. Neural distances compare
realized or prespecified estimated response vectors within a
participant, whereas readout discrimination depends on their conditional
distributions. These are parallel measurements of an occurrence. There
is no assumed causal chain from interpretation to embedding to recorded
neurons, and no assumption that neural coordinates correspond across
participants.

A probe protocol \(\pi\) specifies the questions and their weighting,
answer criterion, norming population, and aggregation rule. Write
\(\rho_i^{(k,\pi)}=P(Y\mid o_i,k,\pi)\) for its independently
characterized outcome distribution. A declared divergence between two
such distributions defines a task-predictive distinction, not a
geometric distance between representations. For example, total variation
gives

\begin{equation}
\Delta_{Y,\pi}(i,j;k)
=\frac12\sum_{y\in\mathcal Y}
\left|\rho_i^{(k,\pi)}(y)-\rho_j^{(k,\pi)}(y)\right|.
\end{equation}

For a neural readout \(g_s\) that returns an outcome distribution,
define its averaged prediction
\(\widehat\rho_i=\mathbb E[g_s(R)\mid o_i,k]\). Relative to the declared
pair and task, an approximate preservation criterion is
\(\operatorname{TV}(\widehat\rho_i,\rho_i)\le\varepsilon\) at each
endpoint. The triangle inequality then gives
\(\operatorname{TV}(\widehat\rho_i,\widehat\rho_j)\ge\Delta_{Y,\pi}(i,j;k)-2\varepsilon\).
This is a sufficient approximate criterion for preserving the declared
outcome distributions, not a necessary-and-sufficient characterization
of all semantic discrimination. When the task-predictive distinction
exceeds twice the error tolerance, the predicted distributions must
remain distinct, without prescribing a Euclidean distance. It must be
evaluated using a declared readout class and held-out observations;
accurate population-average prediction alone need not preserve every
rare distinction. The expectation averages response variability:
matching outcome distributions on average does not establish successful
trial-wise interpretation classification. Section~5.2 gives an
illustrative noise-aware classification analysis, and Section~8
specifies separate held-out tests of distribution preservation and
discrimination.

The substitution score in Section~8 is a separate functional of a
protocol and directional intervention. Write it fully as
\(Q_{ij}^{(k;\pi)}\), with \(Q_{ij}^{(k)}\) used as shorthand only when
\(\pi\) is fixed. A substitution operator must specify which
interpretation is introduced into which receiving context, and the
independently normed comparison must be assigned to occurrence pairs
before neural analysis. The score records probe-relative answer changes
under that intervention; it is not automatically the total variation
between outcomes of two unmodified occurrences. Nor is it confusion
susceptibility, severity, or a universal semantic distance. This
distinction prevents the behavioral score from being inferred from
neural geometry or imported from an item pair without a declared
occurrence assignment.

Functional equivalence is consequently relative to a task and readout
class. Two representation--readout pairs are equivalent for the declared
task when they generate the same conditional outcome predictions over
the eligible domain. Invertible reparameterizations preserve that
equivalence when the transformed readout remains in the allowed class,
but generally change coordinate distances. This observation unifies the
linear-map and CPC examples without asserting equivalence under every
biologically restricted readout. Geometry, predictive information,
recoverability, behavioral consequence, and randomized assignment
effects are distinct evidential objects; the empirical hypothesis links
them rather than identifying them.

\hypertarget{multiple-temporal-horizons}{%
\subsection{Multiple temporal
horizons}\label{multiple-temporal-horizons}}

Lerner et al.~(2011) use scrambled narrative speech to investigate a
hierarchy of cortical temporal receptive windows. Their results motivate
distinguishing sensitivity to local speech structure from integration
over longer narrative intervals. Baldassano et al.~(2017) investigate
event structure in continuous narratives, linking cortical event
representations with hippocampal responses around event boundaries.
These precedents support testing multiscale and event-dependent context,
without establishing a particular decay law for hippocampal semantic
responses.

A candidate analytical representation is
\(c_i\)=(\(h_1\)(\(t_i\)),\ldots,\(h_L\)(\(t_i\)),s(\(e_i\))), where
\(h_l\)(t) integrates features of preceding input under a kernel \(K_l\)
and s represents episode-level information. An exponential kernel is a
parsimonious baseline because it represents recency weighting with one
timescale and permits recursive updating with a finite state; neither
property establishes its biological adequacy. Power-law mixtures and
boundary-gated kernels supply alternatives to compare on held-out
material. For example, an exponential kernel with timescale \(\tau_l\)
supplies a testable baseline; an event-gated kernel can reduce carryover
at an independently annotated episode boundary. These are candidate
predictors, not identified neural mechanisms. Region-specific
combinations of horizons could be investigated along the hippocampal
long axis or in hippocampal--cortical networks, including regions
associated with the default mode network. The model must not assume that
anterior and posterior populations instantiate particular timescales
solely from their anatomy. Comparisons should fit kernels on training
material and test them on held-out episodes, keeping independently
annotated boundaries separate from boundaries estimated using the neural
outcome.

For a feature stream \(\phi(t)\), one such normalized history predictor
is

\begin{equation}
h_l(t)=\int_0^\infty K_l(u)\phi(t-u)\,du,\qquad
K_l(u)=\tau_l^{-1}e^{-u/\tau_l},\quad \tau_l>0.
\end{equation}

For a finite recording, the stream must be defined before its beginning
or the truncated kernel renormalized. An episode gate g(t,u) can set
contributions preceding an independently annotated boundary to zero,
replacing the kernel by \(K_l\)(u)g(t,u) divided by its integral.
Exponential decay encodes a particular recency assumption; abrupt gating
encodes a different assumption about boundary-related updating.
Comparing their held-out performance would test descriptive temporal
models, not identify the biological storage process. Identifiability
also depends on stimulus variation: highly correlated horizons may not
support separate timescale estimates, even if their combined predictor
improves fit.

\hypertarget{event-scripts-task-framing-and-hierarchical-recall}{%
\subsection{Event scripts, task framing, and hierarchical
recall}\label{event-scripts-task-framing-and-hierarchical-recall}}

Research on event scripts motivates enriching context beyond elapsed
time. Baldassano, Hasson, and Norman (2018) identify schematic sequences
shared across restaurant and airport narratives. De Soares et al.~(2024)
examine overlapping location and social scripts and show that top-down
attention shifts behavioral segmentation and the alignment of medial
prefrontal event boundaries. These findings motivate distinguishing
stimulus history from the script made relevant by instructions. They do
not establish that hippocampal word-distance changes arise through the
same mechanism.

An extended analytical context is \(c_i=C(H_{t_i},e_i,S_i,G_i)\), with S
denoting a candidate script and G task framing. Script identity,
temporal horizon, and segmentation are distinct variables: listeners can
hear the same history while weighting different transitions as relevant.
A controlled experiment could hold the recording constant while varying
location-focused or socially focused instructions. Consequence norming
must use the corresponding frame, since relevance depends on the
questions asked. Frame effects may also reflect attention and retrieval
demands, which must remain alternative explanations.

Zhong et al.~(2025) provide a complementary account of narrative
compression. Its random-tree model represents narratives through
hierarchies of summarized segments and constrains retrieval by
working-memory capacity. It addresses quantitative regularities in
recall length and the narrative portions summarized by individual recall
clauses. A model tree node is not an identified hippocampal state or an
information-bottleneck optimum. The account helps distinguish
abstraction during recall from temporal integration during listening and
preservation of task-relevant distinctions. Agreement with aggregate
recall statistics does not determine word-level neural geometry.

\hypertarget{similarity-and-discriminability}{%
\section{Similarity and
discriminability}\label{similarity-and-discriminability}}

An embedding distance can be useful without functioning as a universal
ordering of substitutability. Two expressions may be close because they
concern the same participants, objects, or circumstances while differing
in a relation that comprehension must preserve. This observation
motivates a distinction between broad semantic affinity and contextual
discriminability. The former supports transfer and association; the
latter determines whether a listener can retain a difference required by
the present interpretation.

The reported local reversal is compatible with this distinction. If
closely related meanings would otherwise produce responses that are
difficult to distinguish under biological noise, additional
differentiation could preserve their identity without abolishing their
broader affinity. The argument is functional rather than demonstrative:
it identifies a problem that contrastive coding could solve. The
observational result does not prove that neural responses were moved
apart by an adaptive process, that the relevant words were behaviorally
confused, or that separation improved comprehension.

Several explanations can coexist at this stage. Contextual binding could
create differences that a particular embedding compresses. The measured
distance could change with normalization or representational
dimensionality. Different linguistic properties could contribute
residual structure despite the reported controls. A mechanism supporting
memory differentiation could also produce a semantic pattern without
having been specialized for lexical discrimination. Treating the result
as a question about these alternatives makes the theoretical claim more
useful: it directs inquiry toward the conditions under which local
separation occurs and what that separation accomplishes.

\hypertarget{distance-distortion-and-an-explicit-counterexample}{%
\subsection{Distance distortion and an explicit
counterexample}\label{distance-distortion-and-an-explicit-counterexample}}

The distinction between geometrical correspondence and mechanism can be
made exact. For \(z\in\mathbb R^d\) and \(r\in\mathbb R^n\), let
\(A\in\mathbb R^{n\times d}\) define the linear map \(r=Az\). Squared
Euclidean response distance becomes

\begin{equation}
\|r_i-r_j\|_2^2=(z_i-z_j)^{\mathsf T}A^{\mathsf T}A(z_i-z_j).
\end{equation}

Thus a linear mapping induces a positive-semidefinite quadratic form
\(G=A^{\mathsf T}A\) on embedding differences. If G=aI for a positive
scalar a on the relevant subspace, Euclidean distances are uniformly
rescaled and their ordering is preserved. A general G stretches
directions unequally; distances of two pairs can change order because
their difference vectors point in different directions. This is a
property of the mapping rather than evidence of an added differentiation
process.

For a concrete example, take embedding difference vectors
\(\delta_1\)=(0.1,0) and \(\delta_2\)=(0,0.2), and let A=diag(3,1).
Their original distances are 0.1 and 0.2, but mapped distances are 0.3
and 0.2. The more similar pair becomes more distant in the response
space under a fixed linear transform. This example establishes the
possibility of an ordering change, not a reproduction of the article's
distribution-wide local correlation. Replicating that correlation would
additionally require a specified distribution of pairs and measurement
noise. It nevertheless demonstrates why an observed local ordering
discrepancy cannot identify an adaptive separation mechanism on its own.

For nonzero vectors, cosine distance after mapping is

\begin{equation}
d_{\cos}(Az_i,Az_j)=1-
\frac{z_i^{\mathsf T}Gz_j}
{\sqrt{z_i^{\mathsf T}Gz_i}\sqrt{z_j^{\mathsf T}Gz_j}}.
\end{equation}

The same induced quadratic form changes angles as well as lengths. On
the represented subspace, a positive scalar multiple of an orthogonal
transformation preserves cosine geometry; arbitrary linear maps do not.
Additive baselines and nonlinear transformations can produce further
changes. A consequence-augmented model must therefore be assessed
against mapping baselines of declared complexity, with regularization
and out-of-sample evaluation preventing a flexible mapping from merely
memorizing the tested pairs.

An invertible linear transformation is a homeomorphism of the ambient
Euclidean space, so it preserves the ordinary topology of a represented
support while potentially changing its metric geometry.
Persistent-homology birth and death scales can nevertheless change
because the filtration depends on the selected distances. Conversely,
projection or undersampling can obscure a support's topology. Brown and
Farivar's (2025) analysis consequently motivates an additional testable
description, not a topological explanation of the reported local sign
change. A topology claim here would require a specified support,
sampling assumptions, filtration, and uncertainty analysis; no such
claim follows from the distance correlation alone.

\hypertarget{separation-noise-and-usable-discrimination}{%
\subsection{Separation, noise, and usable
discrimination}\label{separation-noise-and-usable-discrimination}}

Distance has a functional interpretation only under assumptions about
response variability and readout. Consider an illustrative binary
discrimination problem in which response R conditioned on interpretation
a or b follows a multivariate Gaussian with means \(\mu_a\) and
\(\mu_b\) and a common positive-definite covariance Sigma. With equal
class priors, the optimal linear discriminant depends on the squared
Mahalanobis separation

\begin{equation}
\delta^2=(\mu_a-\mu_b)^{\mathsf T}\Sigma^{-1}(\mu_a-\mu_b).
\end{equation}

Under these assumptions, its minimum classification error is
\(\Phi(-\delta/2)\), where \(\Phi\) is the standard normal cumulative
distribution function. To see this, project R onto
\(\Sigma^{-1}(\mu_a-\mu_b)\). The projected class means differ by delta
squared and their common variance is delta squared; a midpoint threshold
therefore lies \(\delta/2\) standard deviations from either mean.
Increasing \(\delta\) reduces error in this model. Increasing raw
Euclidean distance need not do so if the separation grows along a
direction of high or increasing noise. Likewise, two populations with
the same cosine distance can differ in discriminability because their
covariance structures differ.

This calculation is an analytical example, not a Gaussian model fitted
to hippocampal spiking. Neural spike counts may have mean-dependent
variance, correlated fluctuations, and temporal dependence. The
example's purpose is to identify the missing bridge between geometry and
function: a readout-relevant signal-to-noise quantity and an independent
behavioral criterion. A future experiment should evaluate discrimination
using held-out responses or suitably regularized covariance estimation.
Covariance parameters must be estimated without using test observations,
especially when recorded-unit dimension is large relative to repeat
count. Single presentations provide limited information about
condition-specific response distributions; controlled repetitions would
strengthen this part of the functional test.

\hypertarget{predictive-distinction-without-a-prescribed-geometry}{%
\subsection{Predictive distinction without a prescribed
geometry}\label{predictive-distinction-without-a-prescribed-geometry}}

The information bottleneck gives a normative interpretation of relevant
distinction: its distortion compares conditional outcome distributions
rather than points in an embedding or neural space. Under the declared
encoder assumptions, its expected KL distortion is lost predictive
information, or excess expected log loss. The stationary-encoder
condition favors codes whose predictions approximate those required by
the input, but does not require a particular Euclidean or cosine
separation. Appendix A derives this relation and illustrates the
compression trade-off with a binary example. Predictive distortion,
embedding distance, and neural representational distance remain
different mathematical quantities. Likewise, IB code bifurcations do not
identify the reduced-rank predictive dimensions reported by Franch et
al.; relating them would require an additional fitted model.

\hypertarget{operational-hypothesis}{%
\section{Operational hypothesis}\label{operational-hypothesis}}

The central hypothesis is that increasing independently normed,
probe-relative substitution consequence predicts greater neural
representational differentiation beyond lexical, contextual-embedding,
task-difficulty, and fitted-mapping baselines under a controlled
manipulation. Susceptibility to making that substitution is a separate
behavioral quantity, rather than part of the definition of Q. Section~8
specifies how both are measured before neural analysis. An adequately
powered failure to find the predicted differentiation would count
against this functional extension even if the original local sign change
replicated. This is an operational hypothesis, not an optimization model
or a derivation of the reported sign change; no loss function is needed
to state its independent measurement requirement.

For participant s and comparison condition k, let \(r_{si}^{(k)}\) be
the recorded population response to occurrence i, and let \(z_i^{(k)}\)
be its contextual embedding. Define

\begin{equation}
D_{r,sij}^{(k)}=d_r(r_{si}^{(k)},r_{sj}^{(k)}),\qquad
D_{z,ij}^{(k)}=d_z(z_i^{(k)},z_j^{(k)}),\qquad i<j.
\end{equation}

Each neural distance is computed within a participant's population;
neurons from different participants are not assumed to be corresponding
coordinates. Neural and embedding metrics need not be identical, but
their preprocessing, handling of zero vectors, and scaling must be
prespecified. Cosine distance is one candidate for nonzero vectors.
Euclidean distance requires a declared scale and normalization policy.
Let \(Q_{ij}^{(k)}\) denote the symmetric consequence score assigned
from independent norming to an eligible item pair and condition.
Averaging the two directional substitution scores must precede analysis
of a symmetric neural distance; directional scores could instead support
a separate directional behavioral analysis. Only pairs with a defined
normed comparison enter the consequence analysis. Q is not imputed to
every pair in a narrative merely because those occurrences have
embeddings.

Thus \(r_{si}^{(k)}\) belongs to \(\mathbb{R}^{n_s}\), where \(n_s\) is
that participant's recorded unit count; the total of 356 across the
published sample is not the dimension of every participant's response
vector. Cosine distance and the article's word-aligned response window
provide replication-oriented defaults, with alternative metrics and
windows declared as sensitivity analyses. The published 0.15--0.35 BERT
range can define a replication target under matching preprocessing, but
its transfer to a new stimulus set or embedding model requires
independent calibration.

Before choosing a regression, define the associational target. Let \(W\)
collect the declared embedding, linguistic, task, and mapping baselines
for an eligible comparison, with a prespecified participant--stimulus
weighting distribution. Where a differentiable conditional mean is
meaningful, write

\begin{equation}
\mu(q,w)=\mathbb E[D_r\mid Q=q,W=w],\qquad
\tau_Q(q,w)=\frac{\partial\mu(q,w)}{\partial q}.
\end{equation}

The hypothesis predicts a positive conditional consequence association
in the declared regime. This derivative is not a causal effect of
setting the behavioral score. Because the finite probe score has
discrete support, a prespecified difference in conditional means between
supported score levels is an alternative estimand; a smooth derivative
then reflects a working model, not an assumption that every intermediate
score is experimentally realizable. Baselines, support, and target
weighting must be specified before the estimator and test are chosen.

A candidate regression for those eligible pairs is

\begin{equation}
\begin{aligned}
D_{r,sij}^{(k)}={}&\alpha+f(D_{z,ij}^{(k)})
+\theta M_{sij}^{(k)}+\boldsymbol\gamma^{\mathsf T}\mathbf X_{sij}^{(k)}\\
&+(\beta_Q+b_{Qs})Q_{ij}^{(k)}
+\beta_R I_{ij}^{(k)}
+\beta_{QR}Q_{ij}^{(k)}I_{ij}^{(k)}\\
&+b_{0s}+a_i+a_j+v_{si}+v_{sj}+\epsilon_{sij}^{(k)},
\end{aligned}
\end{equation}

where \(I_{ij}^{(k)}\) indicates a prespecified local embedding-distance
range and f is a prespecified smooth baseline, such as a penalized
thin-plate regression spline. M is the distance predicted by a
contextual-embedding-to-neural mapping fitted on independent training
data. The nuisance vector X includes duration contrasts and, where
appropriate, pair means, phonetic distance, declared syntactic features,
frequency, temporal separation, and independently normed susceptibility
and difficulty. A syntactic tree distance must be operationally defined
for occurrences in different sentences rather than presumed to have a
unique value. Smoothness, feature selection, mapping complexity, and
scaling are fixed or selected within training folds. Strong collinearity
among these terms limits the estimability of a unique consequence
association and must be reported.

The participant intercept \(b_{0s}\) and consequence slope \(b_{Qs}\)
may be modeled jointly as mean-zero Gaussian effects with an estimated
covariance. Shared stimulus-occurrence effects \(a_i\) capture item
variation across participants, while participant-specific occurrence
effects \(v_{si}\) capture repeated use of a recorded endpoint in that
participant's distances. The same occurrence effect contributes at
either endpoint, preserving symmetry under exchange of i and j. These
effects are candidate covariance components rather than a guarantee of
independent residuals: distances also share population measurements and
temporal structure. The estimable random-effects structure should be
preregistered subject to declared simplification rules, with
design-based inference and simulation-based calibration used to assess
remaining dependence.

Without the interaction, the proposed conditional association is
\(\beta_Q\)\textgreater0. With the interaction, the outside-range slope
is \(\beta_Q\) and the inside-range slope is \(\beta_Q\)+\(\beta_QR\).
Enhanced local differentiation requires \(\beta_QR\)\textgreater0;
positive differentiation within the local range separately requires
\(\beta_Q\)+\(\beta_QR\)\textgreater0. A positive interaction alone is
insufficient if the total local slope remains negative. Conversely, a
positive local consequence slope does not demonstrate a negative
embedding-distance slope. The latter concerns the derivative of the
embedding-distance relationship, with other predictors controlled, and
must be tested separately. A continuous alternative replaces \(\beta_R\)
I with h(\(D_z\)) and \(\beta_Q\) Q+\(\beta_QR\) QI with Qg(\(D_z\)),
using prespecified smooth functions and simultaneous uncertainty bands.
It avoids hard boundaries but still requires a declared criterion for
stronger consequence effects at high similarity.

Evidence for added predictive information would additionally require
improvement over the same baseline without Q and its interaction on
held-out stimulus blocks. All pairs involving a held-out occurrence must
be excluded from training when testing generalization to new
occurrences. A random split of pairwise rows would leak shared endpoints
across folds. These regressions specify estimands and comparisons; no
coefficients are estimated in this commentary, and an adjusted
consequence association alone would not establish causality.

\hypertarget{identifiability-and-held-out-prediction}{%
\subsection{Identifiability and held-out
prediction}\label{identifiability-and-held-out-prediction}}

The conditional estimand requires overlap: eligible comparisons with
different consequence scores must exist at sufficiently similar baseline
features. If the consequence score is determined by those features, a
partial association cannot be identified by adding more controls. The
regression is one working model for the estimand, rather than its
definition. Nuisance adjustment must address both the outcome and the
consequence predictor, while uncertainty must accommodate shared
endpoints and participants. Appendix B gives the elementary projection
identity and a held-out predictive-gain statistic. The substantive
requirements are a declared direction of association, predictive
improvement on genuinely unseen occurrences, and an independent
discrimination outcome; success on one does not substitute for the
others.

\hypertarget{conditional-association-and-randomized-intervention}{%
\subsection{Conditional association and randomized
intervention}\label{conditional-association-and-randomized-intervention}}

Randomization identifies the effect of the assigned context manipulation
under its design assumptions, not automatically the causal effect of its
derived consequence score. Let \(K\) be a randomized context assignment
and \(D_p(k)\) the potential neural distance for an eligible comparison
under assignment \(k\). Its design-based target is
\(\tau_K=\mathbb E[D_p(1)-D_p(0)]\) under the prespecified comparison
weighting. This assignment effect and \(\tau_Q(q,w)\) are different
estimands and have no general equality. Randomization supports
estimation of this contrast when assignments, interference assumptions,
and missingness are handled appropriately. It does not establish that
every effect of K passes through Q, because context can also change
attention, prediction, acoustic interpretation, or other semantic
variables.

Accordingly, the experiment should separate a primary assignment
contrast from the adjusted Q association. Conditioning on variables
caused by K, such as a context-sensitive surprisal measure, changes the
estimand and may remove part of the total assignment effect. A mediation
interpretation would require further identification assumptions and a
dedicated design. The commentary's principal prediction remains that
consequential distinctions are associated with additional usable
differentiation under controlled conditions, not that a regression
coefficient uniquely measures an adaptive biological mechanism.

\hypertarget{model-alignment-and-predictive-learning-comparisons}{%
\section{Model alignment and predictive-learning
comparisons}\label{model-alignment-and-predictive-learning-comparisons}}

Language-model embeddings provide a tractable comparison space, but
their usefulness should not be confused with their authority as a
definition of meaning. Neural alignment shows that aspects of one
measured organization can be related to aspects of another. It does not
establish shared training objectives, equivalent learning histories,
identical memory operations, or equivalent relationships to action and
experience. Franch et al.~explicitly distinguish shared geometry from
shared algorithm.

The comparison is particularly informative when it reveals disagreement.
A model that captures broad relations while missing a local pattern can
help identify which distinctions its geometry underrepresents relative
to the neural measurement. Yet disagreement alone does not determine
which representation is functionally superior. Superiority requires an
independently specified criterion, such as discrimination performance or
successful contextual interpretation. Likewise, stronger alignment by
one model among those tested establishes a comparative result within
that analysis, rather than a general ranking of architectures as
theories of cognition.

The broader case for vector representations developed by Piantadosi et
al.~(2024) makes such comparisons theoretically intelligible: vectors
can support similarity, relations, and compositional operations within a
common representational framework. The hippocampal findings sharpen a
further question about how these capacities are reconciled when
generalization and differentiation impose competing requirements.
Representational format leaves open the processes that construct and use
it.

Fodor and Suzuki (2026) review 57 fMRI studies using computational
semantic models and identify substantial methodological variation,
alongside robust model--brain associations often dominated by word-level
features. Their synthesis supports asking what measured feature explains
correspondence before interpreting it as a semantic mechanism. Because
their review concerns fMRI rather than hippocampal single-unit geometry,
it supplies a methodological precedent rather than direct evidence for
the present consequence hypothesis.

Cheng, Vaidya, and Antonello (2026) connect intermediate-layer brain
predictivity in language and speech models with intrinsic dimensionality
and semantic abstraction, using fMRI and ECoG comparisons as well as
model interventions. Their results motivate testing representation
properties alongside predictive performance. They do not establish that
prediction is irrelevant to biological comprehension or identify the
origin of the hippocampal local sign change. The distinction is between
an objective that trains a model and the properties of its learned
representation that explain a particular alignment result.

\hypertarget{contrastive-learning-as-a-structural-comparison}{%
\subsection{Contrastive learning as a structural
comparison}\label{contrastive-learning-as-a-structural-comparison}}

Contrastive predictive coding (CPC; van den Oord, Li, and Vinyals, 2018)
trains a context representation to identify a future observation among
sampled alternatives. Under its stipulated sampling scheme, InfoNCE
supplies a lower bound on predictive mutual information; it does not
impose the explicit input-rate penalty of IB. Appendix C states the
objective, density-ratio interpretation, and sampling assumptions. These
frameworks therefore supply comparisons whose relevance to the present
hypothesis can be assessed without assuming that hippocampal learning
instantiates either algorithm.

The key geometrical limitation can be stated directly. CPC's
log-bilinear score is \(f_k(u,c)=\exp\{Z(u)^{\mathsf T}W_kc\}\). For an
invertible matrix \(B\), replacing \(Z(u)\) by \(BZ(u)\) and \(W_k\) by
\(B^{-\mathsf T}W_k\) leaves every predictive score unchanged while
generally changing Euclidean and cosine distances among targets. Thus
functional equivalence does not fix representational geometry, and
improving InfoNCE does not by itself predict the published negative
local distance correlation. The same qualification applies to claims
that increased predictive information entails increased neural distance.

The experimental question is narrower: whether making particular
alternative interpretations available changes held-out discrimination or
neural differentiation at fixed target and independently normed
consequence. Recently activated interpretations, alternative referents,
and recalled episodes are possible sources of competition. A positive
competitor-availability effect would establish functional sensitivity
without identifying biological negative sampling. The behavioral test of
consequential distinction remains intelligible if CPC is rejected as a
model of the underlying neural process.

\hypertarget{experiments-that-distinguish-the-interpretations}{%
\section{Experiments that distinguish the
interpretations}\label{experiments-that-distinguish-the-interpretations}}

A controlled follow-up should vary contextual distinction independently
of lexical similarity. Closely related expressions could appear in
contexts where substituting one for the other either preserves or
changes the interpretation required by a subsequent comprehension
question. Frequency, duration, phonetic similarity, and syntactic
structure would need balancing or explicit modeling. The central outcome
would be whether neural separation tracks the contextual consequence of
substitution beyond what the comparison embedding predicts. This design
would translate the integration--separation question exemplified by
Schlichting et al.~(2015) into a language paradigm, while testing a
different manipulation. Where anatomy and sample size permit, separate
analyses along the hippocampal long axis would help avoid averaging
across potentially distinct coding regimes.

The primary operationalization of the probe-relative substitution
consequence would use an independent behavioral sample, collected and
fixed before the neural data are analyzed. For each item pair and
context condition, norming participants would receive either the
intended interpretation or an explicitly specified substituted
interpretation, followed by the same balanced set of forced-choice
comprehension probes. The answer key would be defined from the narrative
content. Define \(q_{ij}^{(k)}\) as the proportion of probes for which
the normed answer under the substituted interpretation differs from the
answer under the intended interpretation, averaging both substitution
directions where appropriate. Agreement thresholds and exclusion rules
would be prespecified, with poorly understood or inconsistently
interpreted items revised or removed before neural recording. This
measure captures a task-relative consequence of substitution, not
subjective severity or a universal semantic cost. It involves no neural
measurement and does not use embedding distance to define
consequentiality.

The item-pair-to-occurrence assignment used to construct Q must be fixed
before neural analysis. Norming uncertainty should be retained through
resampling of norming participants and probes, or through an explicitly
specified measurement-error model. Probe selection is part of the
estimand: a consequence score depends on which comprehension questions
are sampled. Balancing probe types and testing robustness to alternative
prespecified probe sets would limit dependence on an idiosyncratic
answer key.

Writing the consequence score explicitly clarifies its domain. Let
\(A_{i\to j,l}^{(k)}\) and \(A_{i,l}^{(k)}\) be the independently normed
categorical answers for probe l under a substituted and intended
interpretation, respectively. For L prespecified probes, define

\begin{equation}
q_{i\to j}^{(k)}=\frac{1}{L}\sum_{l=1}^{L}
\mathbf 1\{A_{i\to j,l}^{(k)}\ne A_{i,l}^{(k)}\},\qquad
Q_{ij}^{(k)}=\frac{q_{i\to j}^{(k)}+q_{j\to i}^{(k)}}{2}.
\end{equation}

The symmetric definition requires both directions to be meaningful and
normed under the specified comparison design; otherwise the pair should
be excluded from this analysis or a prespecified alternative estimand
used. A score of one means every sampled answer changes, not that
confusion is maximally harmful in every possible context. Normed
categorical answers should be accompanied by agreement estimates, since
a modal answer can conceal ambiguity. Probability-based
answer-distribution contrasts could provide a separate sensitivity
measure, but their choice and normalization must be declared in advance.
Consequence, confusion probability, and subjective severity are distinct
quantities.

A distinct norming block would present the original speech without
instructing a substitution and measure the rate of foil-consistent
comprehension errors, with reaction time recorded as a secondary index
of difficulty. These outcomes estimate susceptibility and processing
demand separately from q. In the recording sample, the same probe types
would verify comprehension and supply a behavioral outcome; they would
not redefine q after the neural results are known. The primary analysis
would test whether the prespecified consequence measure predicts
held-out neural differentiation after accounting for independently
estimated confusion susceptibility and difficulty. Matching attention
and task demand remains essential, because a consequence effect alone
would not distinguish semantic differentiation from more general task
engagement.

A fuller test of the compression--relevance interpretation would
manipulate relevance separately from the availability of a distinction.
The same recordings could be presented under counterbalanced task
instructions that make a particular semantic contrast informative for
the forthcoming probe in one condition and irrelevant in another.
Independent norming should establish both that the interpretations
remain discriminable under either instruction and that their
contribution to the declared outcome changes. Merely selecting different
probes after seeing neural responses would not constitute this
manipulation. Instructions, probe distributions, incentives, and
analysis contrasts must be fixed in advance, with attention and
difficulty assessed independently. Event-script framing provides a
precedent for altering relevance while preserving the stimulus, but does
not remove these alternative explanations.

For a prespecified matched pair, let \(D_{s,ij}(g)\) denote neural
distance under randomized relevance assignment \(g\in\{0,1\}\). A
participant- and stimulus-aware analysis could test the assignment
contrast

\begin{equation}
\Delta_{\mathrm{rel}}=\mathbb E_{s,(i,j)}
\big[D_{s,ij}(1)-D_{s,ij}(0)\big].
\end{equation}

A positive contrast accompanied by improved held-out recovery of the
task-relevant interpretation would support relevance-sensitive
differentiation under the declared metric and readout. It would not
demonstrate IB optimization: increased distance is only one possible
implementation of improved recovery, and a complete bottleneck test
would also require an identified measure of compression and explicit
model comparison. The narrower consequence experiment tests whether
independently normed substitution consequences explain usable
differentiation; the relevance extension asks whether changing the
task's predictive target changes which distinctions are retained. These
are related hypotheses with distinct estimands, and their randomization
tests must follow their actual assignment units rather than permuting
pairwise cells.

The distribution-preservation criterion of Section 4.1 can be evaluated
in a controlled repeated-condition extension. Estimate the independently
normed outcome distribution for each eligible condition under the same
fixed probe protocol, including probe identity in the outcome when
answers to different questions cannot be pooled meaningfully. Fit a
probabilistic neural readout within training folds, then average its
predictions over held-out responses to estimate \(\widehat\rho_i\).
Endpoint-wise total-variation errors and their participant-, stimulus-,
and norming-aware uncertainty can be compared with a prespecified
tolerance. The criterion provides an interpretable task-distinction
bound only where both endpoint errors are sufficiently small. Trials
supporting response averaging must be repetitions of a declared stimulus
condition or draws from a justified condition ensemble, not silently
pooled distinct narrative occurrences. Separately evaluate trial-wise
interpretation or answer prediction against a baseline that has no
neural response input. Classification accuracy alone does not estimate
total-variation calibration, and this readout extension does not
redefine the substitution score \(Q\) from neural data.

This first experiment tests contextual discriminability. Relational
binding would require a separate manipulation, such as assigning matched
occurrences to the same or different narrative episodes while holding
substitution consequences constant, together with a probe of memory for
that episode association. Such a factorial extension could ask whether
binding modulates discrimination, but it is not necessary to the
narrower claim and is not inferred from distance measurements alone.

A complementary experiment could hold an ambiguous word constant while
varying cues that establish different interpretations. Such a design
would separate lexical identity from contextual differentiation and
allow repeated presentations within each condition. Independent
behavioral judgments would help distinguish successful disambiguation
from neural variability without a demonstrated semantic consequence.
Recording across the relevant sequence could also clarify whether
separation follows the disambiguating cue, precedes it through
expectation, or emerges through later integration.

Acoustic and lexical baselines should be specified at the occurrence
level before pairwise contrasts are constructed. Candidate controls
include intensity, pitch and prosodic prominence, phoneme identity and
timing, coarticulation, lexical frequency, cohort size, phonological
neighborhood density, and contextual surprisal, alongside duration and
syntactic structure. Frequency is a lexical distributional measure,
whereas surprisal is conditional on a specified predictive model and
context; they cannot be substituted for one another. Closely adjacent
response windows may also share acoustic or neural activity, requiring
temporal-separation controls and declared overlap exclusions. Where
feasible, identical target recordings embedded in counterbalanced
contexts would reduce target-acoustic variation, while preserving
natural prosody would require explicit checks rather than assuming
splicing is innocuous.

These controls do not make semantics statistically pure. Semantic,
syntactic, and predictive features covary, and indiscriminately
orthogonalizing embeddings against acoustic or surprisal predictors can
remove task-relevant structure or make conclusions depend on
residualization order. Baseline and consequence-augmented models should
instead be compared on the same held-out data; any residualization must
be fitted within training folds, with its direction declared. Unique
predictive improvement is conditional on the included baselines, not an
exhaustive attribution of residual variance to semantic distinction.
Because surprisal and task engagement may mediate effects of context,
analyses with and without them should distinguish a conditional
association from the total effect of a randomized context manipulation.

Analytical robustness is equally necessary. Similarity ranges and
distance measures should be prespecified using an independent stimulus
set or selected in a training partition and evaluated in held-out
narratives. Report the full association alongside local slopes and their
uncertainty, including sensitivity to shifted range endpoints and a
continuous nonlinear model that does not require a hard boundary.
Repetition across BERT and GPT-2 is informative, but ranges should be
calibrated separately for each representation rather than assuming that
equal numerical distances imply equal semantic similarity. A sign change
that depends narrowly on one endpoint choice would weaken the
interpretation of a distinct coding regime.

Pairwise comparisons reuse occurrences and cannot be treated as millions
of independent biological replications. One concrete approach is to
estimate participant-specific slopes or contrasts and summarize them
with participant-level uncertainty. Another is a multilevel model with
participant effects and crossed occurrence effects for both members of a
pair, supplemented by stimulus or narrative effects where the design
permits. Label permutations must act on the occurrence-level assignment
and regenerate the full distance matrix, preserving shared endpoints,
rather than independently shuffling pairwise cells. For naturalistic
speech, the permutation scheme should respect exchangeability through
matched blocks or a justified temporal procedure; arbitrary shuffling
can destroy narrative dependencies. Hierarchical bootstrap resampling of
participants and narrative blocks provides a complementary sensitivity
analysis. These recommendations specify alternatives for future testing
and do not imply that the published analyses omitted all such controls.

The permutation null must match the regression claim. Shuffling
occurrence labels across a whole narrative tests a broad alignment null
and generally disrupts the legitimate relations among Q, embeddings, and
nuisance features; it is not automatically a valid test of the adjusted
Q coefficient. In a randomized experiment, the primary randomization
test would reassign the context or consequence manipulation at its
actual assignment unit within the permitted blocks. For every
reassignment, one must reconstruct affected occurrence labels and
pair-level predictors and refit the same preregistered analysis,
including any data-dependent steps permitted within training folds. A
whole-matrix label permutation, when exchangeable, must use the same
permutation for rows and columns rather than permuting cells
independently. If several participants receive the same randomized item
assignment, the test must preserve that shared assignment structure
rather than independently permuting each participant.

Where Q varies observationally and no defensible assignment or
exchangeability scheme exists, an exact permutation test of its
conditional effect should not be claimed. Participant- and
stimulus-aware uncertainty estimation, a calibrated model-based
bootstrap, and held-out prediction are alternatives with their own
assumptions. Stimulus blocks shared across participants should be
resampled jointly in a crossed participant--stimulus procedure,
preserving endpoint identities and reconstructing distances, rather than
resampled independently within each participant. With few narratives,
stimulus-generalization uncertainty may remain poorly determined. Null
simulations should assess type-I error under shared-endpoint noise and
temporal dependence, and the primary coefficient or contrast,
local-range rule, multiplicity adjustment, and smallest effect of
interest should be fixed before testing. An inconclusive interval is not
evidence against the hypothesis; a sufficiently precise interval
excluding the prespecified meaningful positive effect would provide the
relevant negative result.

A reproduction across independently specified ranges and measures would
strengthen the geometrical claim, while a relation between
experimentally manipulated confusion consequences and held-out neural
differentiation would strengthen the proposed functional extension.
Learning-related change or intervention would be needed to identify the
mechanism. This separates three possible achievements: reproducing a
sign change, explaining its relation to comprehension, and establishing
the process that produces it.

\hypertarget{conclusion}{%
\section{Conclusion}\label{conclusion}}

The hippocampal study offers evidence that semantic organization during
narrative listening is distributed, contextual, and not uniformly
captured by a simple correspondence between embedding distance and
neural distance. Its theoretical interest lies in the conjunction of
broad relational structure with local differentiation. A listener must
preserve the affinities through which expressions become intelligible
and the distinctions through which their present interpretation remains
determinate.

The contextual-discrimination hypothesis proposed here treats those
requirements as complementary constraints on representation, with
relational memory providing motivation rather than an inferred binding
mechanism. Context changes which similarities can support generalization
and which differences must remain recoverable. The reported local sign
change makes this proposal worth testing, while leaving its adaptive and
mechanistic explanations open. A mature account of semantic geometry
would therefore connect population organization to independently
measured discrimination and comprehension, explaining how continuity
across contexts can coexist with the preservation of consequential
differences.

\appendix

\hypertarget{predictive-distortion-and-the-information-bottleneck}{%
\section{Predictive distortion and the information
bottleneck}\label{predictive-distortion-and-the-information-bottleneck}}

The information-bottleneck framework permits an exact statement of why
similar inputs may need distinct predictive codes. Let \(p(x,y)\) be a
joint distribution, and let a stochastic encoder \(p(t\mid x)\) satisfy
the Markov condition \(Y-X-T\), meaning that the representation is
generated from X without additional access to Y. For discrete variables
define

\begin{equation}
d_{\mathrm{IB}}(x,t)=D_{\mathrm{KL}}\bigl(p(Y|x)\|p(Y|t)\bigr),\qquad
D_{\mathrm{KL}}(p\|q)=\sum_y p(y)\log\frac{p(y)}{q(y)}.
\end{equation}

The divergence is directed and can be infinite when q assigns zero
probability to an outcome having positive probability under p.~It is not
a metric between neural vectors. Its expectation satisfies

\begin{equation}
\mathbb E_{X,T}d_{\mathrm{IB}}(X,T)
=H(Y|T)-H(Y|X)=I(X;Y)-I(T;Y)\geq0.
\end{equation}

Expanding the expected logarithm and using conditional independence
gives the first equality; the definition of mutual information gives the
second. The quantity measures lost predictive information, equivalently
excess expected logarithmic prediction loss, under the stated model. It
cannot be equated with cosine distance or applied to arbitrary neural
observations without a probabilistic specification.

Tishby et al.'s stationary encoder equation makes the role of distortion
explicit:

\begin{equation}
p(t|x)=\frac{p(t)}{Z_\beta(x)}e^{-\beta d_{\mathrm{IB}}(x,t)},\qquad
Z_\beta(x)=\sum_t p(t)e^{-\beta d_{\mathrm{IB}}(x,t)}.
\end{equation}

The representative prediction \(p(y\mid t)\) depends on the encoder and
joint distribution, so this is a self-consistency condition. At fixed
representative predictions, assigning an input to a code with greater
predictive distortion receives less probability, other factors equal.
Divergent predictions can therefore motivate retaining different codes,
but they do not require those codes to have greater neural Euclidean or
cosine distance.

For an elementary example, let X take equally probable values a and b,
with \(p(Y=1\mid a)=\epsilon\) and \(p(Y=1\mid b)=1-\epsilon\) for
\(0<\epsilon<1/2\). Merging them into a single code gives
\(p(Y=1\mid t)=1/2\) and loses predictive information

\begin{equation}
\log2-H_{\mathrm b}(\epsilon),\qquad
H_{\mathrm b}(\epsilon)=-\epsilon\log\epsilon
-(1-\epsilon)\log(1-\epsilon).
\end{equation}

A deterministic two-code representation removes that loss while raising
\(I(X;T)\) from zero to log 2. Comparing only these two candidate
encoders, retention is favored when \(\beta\) times the lost predictive
information exceeds log 2. Stochastic encoders can produce intermediate
trade-offs, so the comparison is not a general phase-transition result.
It illustrates why arbitrarily close comparison embeddings can conceal
consequential predictive differences, without deriving any
neural-distance sign change.

Tishby et al.~also discuss bifurcations of code solutions as \(\beta\)
changes in the information plane. Such changes concern predictive
quantizations and code cardinality. Franch et al.'s reduced-rank
regression instead estimates linear predictive dimensions for different
embedding models. Rank 12 for GPT-2 and rank 3 for Word2Vec therefore do
not demonstrate bottleneck bifurcations or hierarchical code splitting.
Connecting the quantities would require a fitted generative model
relating code structure to regression rank.

Futrell and Hahn (2026) offer a related but distinct
sequential-processing account. They define predictive information as
dependence between a process's past and future and show that codes
constrained to have low predictive information can express approximately
independent meaning features systematically and locally, with
cross-linguistic corpus comparisons supporting their structural
argument. This concerns the organization of linguistic signals under
processing constraints. It is not evidence that a hippocampal
representation minimizes the IB objective, and minimizing predictive
dependence within a signal is not equivalent to maximizing CPC's
predictive-information bound. The contribution is a precedent for
deriving structural hypotheses from explicitly specified information
constraints rather than inferring an objective from alignment alone.

\hypertarget{adjustment-and-predictive-improvement}{%
\section{Adjustment and predictive
improvement}\label{adjustment-and-predictive-improvement}}

In an ordinary linear regression with independent errors, the
conditional coefficient can be written using projections. For \(N\)
eligible observations, let \(B\in\mathbb R^{N\times p}\) be the baseline
design matrix and \(y,q\in\mathbb R^N\) the outcome and consequence
vectors. Write \(P_B=BB^+\) for the orthogonal projection onto the
column space of \(B\), where \(B^+\) is its Moore--Penrose
pseudoinverse, and set \(M_B=I_N-P_B\). Both projection operators are
\(N\times N\); when \(B\) has full column rank,
\(P_B=B(B^{\mathsf T}B)^{-1}B^{\mathsf T}\). Provided
\(q^{\mathsf T}M_Bq\) is positive,

\begin{equation}
\widehat\beta_Q=
\frac{q^{\mathsf T}M_B y}{q^{\mathsf T}M_B q}.
\end{equation}

This expression shows why nuisance adjustment applies to both the
outcome and the consequence predictor: it compares their components not
accounted for by B. Residualizing y alone and regressing it on
unadjusted q is generally a different procedure. The expression is
explanatory; the proposed hierarchical analysis requires
covariance-aware estimation and cannot import independent-error standard
errors from this elementary case. Small residual variation in q makes
estimation unstable even when its raw variation is substantial.

Predictive improvement can be stated independently of coefficient
significance. For held-out eligible pairs \(P_test\), define

\begin{equation}
\Delta_{\mathrm{pred}}=
\frac{1}{|P_{\mathrm{test}}|}\sum_{p\in P_{\mathrm{test}}}
\left[(y_p-\widehat y_{0p})^2-(y_p-\widehat y_{1p})^2\right],
\end{equation}

where model 0 contains the declared baselines and model 1 additionally
includes consequence and the prespecified interaction. Positive
improvement indicates lower test error for the augmented model. Pairwise
test losses remain dependent, so uncertainty must respect their
participant and stimulus structure. Weighting should also be specified:
equal weighting of pairs favors participants or narratives contributing
many eligible comparisons, whereas averaging participant-specific losses
targets a different population summary. Coefficient direction,
predictive gain, and behavioral discrimination are complementary
criteria rather than interchangeable evidence.

\hypertarget{contrastive-predictive-coding-objective-and-assumptions}{%
\section{Contrastive predictive coding: objective and
assumptions}\label{contrastive-predictive-coding-objective-and-assumptions}}

Contrastive predictive coding (CPC; van den Oord, Li, and Vinyals, 2018)
offers a predictive-learning comparison whose assumptions can be stated
explicitly. Its encoder constructs latent observations, while an
autoregressive network summarizes their history into a context
representation. To avoid confusing the CPC notation with the
commentary's situated context, write this learned summary as \(C_t\) and
a future observation as \(U=x_{t+k}\). For each context, a candidate set
contains one positive observation drawn from \(p(u\mid C_t)\) and
\(N-1\) independent alternatives drawn from the marginal \(p(u)\). With
a positive scoring function \(f_k\), its population objective and
information bound are

\begin{equation}
\mathcal L_{N,k}=-\mathbb E\log
\frac{f_k(U^+,C_t)}{\sum_{j=1}^{N}f_k(U_j,C_t)},\qquad
I(U;C_t)\ge\log N-\mathcal L_{N,k}.
\end{equation}

The expectation is over the joint context--positive distribution and the
stipulated negative-sampling distribution. The bound concerns that
population objective, rather than every finite-sample estimate, and its
maximum possible value is \(\log N\). It is a lower bound on predictive
information, not an estimate of the information rate compressed from the
entire speech history. In particular, CPC does not include the explicit
input-rate penalty in the IB objective. Architectural restrictions may
constrain capacity, but a compact vector and a penalized
mutual-information rate are different commitments. Thus IB and CPC
supply related perspectives on predictive retention without constituting
equivalent optimization problems.

The density-ratio interpretation follows from the
candidate-identification problem. When the positive index is uniformly
distributed, the Bayes posterior is proportional to
\(p(U_j\mid C_t)/p(U_j)\). An unrestricted optimal score consequently
recovers this ratio up to a context-dependent multiplicative factor. CPC
uses a log-bilinear critic; writing its target latent vector as \(Z(u)\)
gives

\begin{equation}
f_k(u,c)\propto\frac{p(u\mid c)}{p(u)},\qquad
f_k(u,c)=\exp\{Z(u)^{\mathsf T}W_kc\}.
\end{equation}

The first relation describes the ideal density-ratio solution, whereas
the second specifies a restricted parameterization that need not attain
it. Neither relation is a metric between two word-response vectors.
Indeed, replacing \(Z(u)\) by \(BZ(u)\) for an invertible matrix \(B\)
and \(W_k\) by \(B^{-\mathsf T}W_k\) leaves every score unchanged while
potentially altering Euclidean and cosine distances among target
vectors. This invariance parallels the mapping concern in Section~5:
equivalent predictive scores do not determine a unique distance
geometry. A competitor can receive a lower score through a changed
critic or context representation without requiring increased pairwise
neural distance. Under these assumptions, optimizing InfoNCE therefore
does not predict a negative local semantic--neural distance correlation.

The selection of alternatives is equally consequential. The standard
bound assumes marginal negatives. Sampling only semantically similar
alternatives, or excluding same-episode observations, changes the
candidate distribution and requires a correspondingly derived objective;
the original mutual-information interpretation cannot simply be carried
over. False negatives are also possible when two observations preserve
the same relevant interpretation. CPC predicts sampled future
observations, whereas the commentary's behavioral variable concerns
consequences for declared comprehension probes. Future predictability
can retain acoustic regularities or other features without preserving
the intended task distinction, so identifying these targets would
require evidence rather than terminology. BERT and GPT-2 in the
hippocampal comparison are not thereby shown to have been trained with
InfoNCE.

The biological question is what would select the alternatives. Candidate
competitors include recently activated interpretations, alternative
referents, or recalled episodes sharing cues. Experimentally
manipulating which alternatives are available, while holding the target
and its independently normed consequence score constant, could test
whether differentiation depends on competitor availability. Such an
effect would support a functional role for competition without
identifying negative sampling, a density-ratio critic, or InfoNCE
optimization in neural circuitry. The joint information-theoretic
interpretation is therefore deliberately conditional: IB specifies which
declared predictive distinctions compression should preserve, and CPC
illustrates one way to train predictive scores against sampled
alternatives. A model that connects either objective to the reported
hippocampal geometry must additionally specify its representation,
noise, readout, and competitor distribution.
\clearpage
\begin{thebibliography}{99}
\bibitem{ref1}
Bakker, A., Kirwan, C. B., Miller, M., and Stark, C. E. L. (2008).
Pattern separation in the human hippocampal CA3 and dentate gyrus.
\emph{Science}, 319(5870), 1640--1642.
https://doi.org/10.1126/science.1152882.

\bibitem{ref2}
Baldassano, C., Chen, J., Zadbood, A., Pillow, J. W., Hasson, U., and
Norman, K. A. (2017). Discovering event structure in continuous
narrative perception and memory. \emph{Neuron}, 95(3), 709--721.e5.
https://doi.org/10.1016/j.neuron.2017.06.041.

\bibitem{ref3}
Baldassano, C., Hasson, U., and Norman, K. A. (2018). Representation of
real-world event schemas during narrative perception. \emph{Journal of
Neuroscience}, 38(45), 9689--9699.
https://doi.org/10.1523/JNEUROSCI.0251-18.2018.

\bibitem{ref4}
Brown, S., and Farivar, R. (2025). The topology of representational
geometry. \emph{Frontiers in Neuroscience}, 19, 1597899.
https://doi.org/10.3389/fnins.2025.1597899.

\bibitem{ref5}
Cheng, E., Vaidya, A. R., and Antonello, R. (2026). Abstraction induces
the brain alignment of language and speech models. In \emph{Proceedings
of the 43rd International Conference on Machine Learning}, PMLR 306,
19008--19031. https://proceedings.mlr.press/v306/cheng26l.html.

\bibitem{ref6}
De Soares, A., Kim, T., Mugisho, F., Zhu, E., Lin, A., Zheng, C., and
Baldassano, C. (2024). Top-down attention shifts behavioral and neural
event boundaries in narratives with overlapping event scripts.
\emph{Current Biology}, 34(20), 4729--4742.e5.
https://doi.org/10.1016/j.cub.2024.09.013.

\bibitem{ref7}
Deniz, F., Tseng, C., Wehbe, L., Dupré la Tour, T., and Gallant, J. L.
(2023). Semantic representations during language comprehension are
affected by context. \emph{Journal of Neuroscience}, 43(17), 3144--3158.
https://doi.org/10.1523/JNEUROSCI.2459-21.2023.

\bibitem{ref8}
Duff, M. C., and Brown-Schmidt, S. (2012). The hippocampus and the
flexible use and processing of language. \emph{Frontiers in Human
Neuroscience}, 6, 69. https://doi.org/10.3389/fnhum.2012.00069.

\bibitem{ref9}
Fodor, J., and Suzuki, S. (2026). Using computational semantics to study
meaning in the brain. \emph{Neuroscience \& Biobehavioral Reviews}, 181,
106514. https://doi.org/10.1016/j.neubiorev.2025.106514.

\bibitem{ref10}
Franch, M., Mickiewicz, E. A., Belanger, J. L., et al.~(2026). A
population code for semantics in human hippocampus. \emph{Nature
Neuroscience}. https://doi.org/10.1038/s41593-026-02436-4.

\bibitem{ref11}
Futrell, R., and Hahn, M. (2026). Linguistic structure from a bottleneck
on sequential information processing. \emph{Nature Human Behaviour}, 10,
589--600. https://doi.org/10.1038/s41562-025-02336-w.

\bibitem{ref12}
Goldstein, A., Grinstein-Dabush, A., Schain, M., et al.~(2024).
Alignment of brain embeddings and artificial contextual embeddings in
natural language points to common geometric patterns. \emph{Nature
Communications}, 15, 2768. https://doi.org/10.1038/s41467-024-46631-y.

\bibitem{ref13}
Goldstein, A., Zada, Z., Buchnik, E., et al.~(2022). Shared
computational principles for language processing in humans and deep
language models. \emph{Nature Neuroscience}, 25, 369--380.
https://doi.org/10.1038/s41593-022-01026-4.

\bibitem{ref14}
Konkel, A., and Cohen, N. J. (2009). Relational memory and the
hippocampus: Representations and methods. \emph{Frontiers in
Neuroscience}, 3(2), 166--174.
https://doi.org/10.3389/neuro.01.023.2009.

\bibitem{ref15}
Kriegeskorte, N., Mur, M., and Bandettini, P. (2008). Representational
similarity analysis: Connecting the branches of systems neuroscience.
\emph{Frontiers in Systems Neuroscience}, 2, 4.
https://doi.org/10.3389/neuro.06.004.2008.

\bibitem{ref16}
Lerner, Y., Honey, C. J., Silbert, L. J., and Hasson, U. (2011).
Topographic mapping of a hierarchy of temporal receptive windows using a
narrated story. \emph{Journal of Neuroscience}, 31(8), 2906--2915.
https://doi.org/10.1523/JNEUROSCI.3684-10.2011.

\bibitem{ref17}
McClelland, J. L., McNaughton, B. L., and O'Reilly, R. C. (1995). Why
there are complementary learning systems in the hippocampus and
neocortex: Insights from the successes and failures of connectionist
models of learning and memory. \emph{Psychological Review}, 102(3),
419--457. https://doi.org/10.1037/0033-295X.102.3.419.

\bibitem{ref18}
Patterson, K., Nestor, P. J., and Rogers, T. T. (2007). Where do you
know what you know? The representation of semantic knowledge in the
human brain. \emph{Nature Reviews Neuroscience}, 8, 976--987.
https://doi.org/10.1038/nrn2277.

\bibitem{ref19}
Piantadosi, S. T., Muller, D. C. Y., Rule, J. S., Kaushik, K.,
Gorenstein, M., Leib, E. R., and Sanford, E. (2024). Why concepts are
(probably) vectors. \emph{Trends in Cognitive Sciences}, 28(9),
844--856. https://doi.org/10.1016/j.tics.2024.06.011.

\bibitem{ref20}
Rigotti, M., Barak, O., Warden, M. R., Wang, X.-J., Daw, N. D., Miller,
E. K., and Fusi, S. (2013). The importance of mixed selectivity in
complex cognitive tasks. \emph{Nature}, 497, 585--590.
https://doi.org/10.1038/nature12160.

\bibitem{ref21}
Schlichting, M. L., Mumford, J. A., and Preston, A. R. (2015).
Learning-related representational changes reveal dissociable integration
and separation signatures in the hippocampus and prefrontal cortex.
\emph{Nature Communications}, 6, 8151.
https://doi.org/10.1038/ncomms9151.

\bibitem{ref22}
Tishby, N., Pereira, F. C., and Bialek, W. (1999; arXiv deposit 2000).
The information bottleneck method. \emph{Proceedings of the 37th Annual
Allerton Conference on Communication, Control, and Computing}. Author
manuscript: https://arxiv.org/abs/physics/0004057.

\bibitem{ref23}
van den Oord, A., Li, Y., and Vinyals, O. (2018; revised 2019).
Representation learning with contrastive predictive coding. \emph{arXiv
preprint}, arXiv:1807.03748. https://arxiv.org/abs/1807.03748.

\bibitem{ref24}
Yan, X., Chavez, A. G., Krishna, A., et al.~(2026). Shared neural
geometries for bilingual semantic representations in human hippocampal
neurons. \emph{Cell}, 189(16), 5065--5080.e10.
https://doi.org/10.1016/j.cell.2026.05.020.

\bibitem{ref25}
Zhong, W., Can, T., Georgiou, A., Shnayderman, I., Katkov, M., and
Tsodyks, M. (2025). Random tree model of meaningful memory.
\emph{Physical Review Letters}, 134, 237402.
https://doi.org/10.1103/g1cz-wk1l.

\end{thebibliography}
\end{document}

